% Демонстрация 5: «Комплексные числа и корни n-й степени» (Лекция 12)
\section*{Демонстрация 5. «Комплексные числа и корни \texorpdfstring{$n$}{n}-й степени»\\ \normalsize\textmd{Файл \texttt{demos/05-complex-numbers.html} \quad(к Лекции~12)}}

\subsection*{Какие факты лекции иллюстрирует}

Интерактив наглядно связывает геометрию комплексной плоскости с
формулами лекции. Перетаскивая точку $z$ мышью или касанием,
преподаватель в реальном времени демонстрирует следующие положения.

\begin{itemize}
  \item \textbf{Геометрический смысл модуля и аргумента.} Радиус-вектор
  точки $z = x + iy$ показывает длину $r = |z| = \sqrt{x^2 + y^2}$ и
  угол наклона $\varphi = \arg z$ к действительной оси; панель
  одновременно выводит алгебраическую запись, модуль, аргумент (в
  радианах и градусах) и дугу угла $\varphi$ на плоскости. Это оживляет
  переход от декартовых координат к полярным.

  \item \textbf{Тригонометрическая и показательная формы.} Для текущей
  точки сразу показаны записи
  \[
    z = r(\cos\varphi + i\sin\varphi) = r e^{i\varphi},
  \]
  так что студент видит: это \emph{одно и то же число} в трёх обличьях.

  \item \textbf{Умножение как поворот и растяжение.} В режиме
  «Степень~$z^n$» при $n = 2$ хорошо видно, что $z^2$ получается
  удвоением аргумента и возведением модуля в квадрат — частный случай
  правила «при умножении аргументы складываются, модули
  перемножаются».

  \item \textbf{Формула Муавра.} Увеличивая степень ползунком $n$,
  наблюдаем, как точка $z^n$ движется по спирали; панель выводит
  значение и запись
  \[
    z^n = r^n(\cos n\varphi + i\sin n\varphi).
  \]
  Если поставить $z$ на единичную окружность ($r = 1$), то $z^n$ остаётся
  на ней, и возведение в степень превращается в чистое вращение на угол
  $n\varphi$.

  \item \textbf{Корни $n$-й степени как вершины правильного
  многоугольника.} В режиме «Корни~$\sqrt[n]{z}$» демонстрация рисует
  все $n$ корней на окружности радиуса $r^{1/n}$ и соединяет их в
  правильный $n$-угольник, отражая формулу
  \[
    \sqrt[n]{z} = r^{1/n}\!\left(
      \cos\frac{\varphi + 2\pi k}{n} + i\sin\frac{\varphi + 2\pi k}{n}
    \right), \quad k = 0, 1, \dots, n-1.
  \]
\end{itemize}

\begin{remarkIT}{Где это работает в технике}{demo05_apps}
  Симметрия корней из единицы, которую демонстрирует интерактив, — не
  только красивая картинка. Равномерно расставленные по окружности корни
  $e^{-i 2\pi k/N}$ лежат в основе \textbf{дискретного преобразования
  Фурье} и его быстрого алгоритма (ДПФ/БПФ), а значит, кодеков MP3 и
  JPEG, OFDM-модуляции в Wi-Fi. Умножение как поворот применяется в
  \textbf{2D-графике и геймдеве} (поворот спрайта — одно комплексное
  умножение), а в электротехнике комплексные числа-\textbf{фазоры}
  описывают амплитуду и фазу гармонических токов. Наконец, итерация
  $z \mapsto z^2 + c$ на комплексной плоскости порождает
  \textbf{фракталы} (множества Мандельброта и Жюлиа) — пример
  комплексной динамики.
\end{remarkIT}

\subsection*{Примерный сценарий использования на занятии}

Демонстрацию удобно включать после введения тригонометрической формы и
держать открытой до конца разговора о корнях. Ниже — опорный сценарий.

\begin{enumerate}
  \item \textbf{Запуск (после определения модуля и аргумента).} Откройте
  интерактив, перетащите $z$ в первую четверть и покажите соответствие:
  длина вектора — это $|z|$, дуга — это $\varphi$. Подведите точку к
  значениям $i$, $-1$, $1+i$ и сверьте показания панели с устным счётом.

  \item \textbf{Три формы — одно число.} Не двигая точку, обратите
  внимание класса на то, что алгебраическая, тригонометрическая и
  показательная записи в панели описывают одну и ту же точку. Спросите:
  «какая форма удобнее для сложения, а какая — для умножения?»

  \item \textbf{Возведение в степень = вращение.} Переведите $z$ строго
  на единичную окружность ($r = 1$) и включите режим «Степень». Двигая
  ползунок $n$, покажите, что точка $z^n$ бежит по той же окружности,
  поворачиваясь на $n\varphi$. Вопрос к аудитории: «почему модуль не
  меняется?» — потому что $1^n = 1$.

  \item \textbf{Спираль для $r \neq 1$.} Сдвиньте $z$ за пределы
  окружности и снова прокрутите $n$: теперь точка уходит по спирали,
  ведь $r^n$ растёт. Свяжите наблюдаемое с формулой Муавра
  $z^n = r^n(\cos n\varphi + i\sin n\varphi)$ на доске.

  \item \textbf{Постановка вопроса о корнях.} Переключитесь в режим
  «Корни» и задайте вопрос: «Сколько корней у уравнения $z^n = 1$ и где
  они расположены?» Дайте классу высказать гипотезы, прежде чем
  показывать ответ.

  \item \textbf{Момент удивления.} Поставьте $z = 1$, $n = 5$ — на экране
  возникает правильный пятиугольник из пяти корней. Меняя $n$ ползунком,
  покажите, что корни \emph{всегда} образуют правильный $n$-угольник,
  вписанный в окружность радиуса $r^{1/n}$, как бы ни лежала точка $z$.
  Это и есть ключевой визуальный эффект демонстрации.

  \item \textbf{Связь с формулой и выводом.} Перетащите $z$ в
  произвольное место и обратите внимание: поворот всего «веера» корней
  отслеживает аргумент $\varphi$, а вершины отстоят друг от друга на
  $2\pi/n$. Отсюда естественно выписывается формула корней
  $\dfrac{\varphi + 2\pi k}{n}$, $k = 0, \dots, n-1$, и становится
  понятно, почему различных значений ровно $n$ (далее углы повторяются с
  периодом $2\pi$).

  \item \textbf{Закрепление на примере лекции.} Воспроизведите в
  интерактиве разобранный в лекции случай $\sqrt[3]{-1}$: поставьте
  $z = -1$, $n = 3$ и сверьте три вершины треугольника с найденными
  аналитически корнями $\tfrac12 \pm i\tfrac{\sqrt3}{2}$ и $-1$. Так
  студент видит совпадение «картинки» и «вычисления».
\end{enumerate}
